\documentclass{article}
\usepackage{graphicx} % Required for inserting images
\usepackage{xeCJK} %
\usepackage{amsmath} %

\title{國中基礎理化公式整理}
\author{By1437大帝}
\date{26/2/26}

\begin{document}

\maketitle

\section{公式}
\[\begin{aligned}
&\text{第一單元：} M=V \cdot D \\
&\text{第二單元：} \text{溶解度}=\frac{\text{溶質}}{100g \text{水}} \text{ , } \text{濃度}=\frac{\text{溶質}}{\text{溶質}+\text{溶劑}} \\
&\qquad\qquad 1 \text{ ppm} = 1 \text{ mg/kg} = 1 \text{ mg/L} \\
&\text{第三單元：} V=331+0.6t \text{ , } V=f \times \lambda \\
&\text{第四單元：} V=f \cdot \lambda \\
&\text{第五單元：} H=ms\Delta T \\
&\text{第六單元：} \text{質量數}=\text{質子數}+\text{中子數} \text{ , } \text{原子序}=\text{質子數}=\text{電子數} \\
&\text{第七單元：} \text{莫耳數}=\frac{\text{質量}}{\text{分子量}}=\frac{\text{粒子數}}{6 \times 10^{23}} \\
&\text{第八單元：} M = \frac{\text{溶質莫耳數(mole)}}{\text{溶液體積(L)}} \\
&\qquad\qquad [H^+] \times [OH^-]=10^{-14}M \text{ (25℃)} \\
&\qquad\qquad pH=-\log[H^+] \text{ , } pOH=-\log[OH^-] \\
&\text{第九單元：} E_k=\frac{1}{2}MV^2=\frac{3}{2}kT \\
&\text{第十一單元：} F = KX \text{ , } (f_s)_{max} = \mu_s N \text{ , } P = \frac{F}{A} \\
&\qquad\qquad P=dgh \text{ (or } hd \text{)} \\
&\qquad\qquad 1 \text{ atm} = 76 \text{ cmHg} = 1033.6 \text{ gw/cm}^2 = 1013.25 \text{ hPa} \\
&\qquad\qquad B_{\text{浮}} = V_{\text{下}} \cdot D_{\text{液}} \text{ , } B_{\text{沉}} = V_{\text{物}} \cdot D_{\text{液}} \\
&\text{第十二單元：} a=\frac{\Delta V}{\Delta t} \text{ , } V_t=V_0+at \text{ , } \Delta X=V_0t+\frac{1}{2}at^2 \\
&\qquad\qquad V^2=V_0^2+2a\Delta X \text{ , } V=\frac{\mathrm{d}X}{\mathrm{d}t} \text{ , } a=\frac{\mathrm{d}V}{\mathrm{d}t} \\
&\text{第十三單元：} F=ma \text{ , } F=G \frac{m_1 m_2}{r^2} \\
&\text{第十四單元：} W=FS\cos\theta \text{ , } E_k=\frac{1}{2}MV^2 \text{ , } E_p=mgh \\
&\qquad\qquad \tau=F \cdot d \text{ , } F=W \cdot \sin\theta \\
&\text{第十五單元：} F=K \frac{Q_1 Q_2}{r^2} \text{ , } I=\frac{Q}{t} \text{ , } V=\frac{E}{Q} \text{ , } R=\frac{V}{I} \text{ , } R=\rho \frac{L}{A} \\
&\text{第十六單元：} E=QV \text{ , } P=\frac{E}{t} \text{ , } 1\text{度}=3.6 \times 10^6 J \\
&\text{第十七單元：} B=\frac{\mu_0 I}{2\pi r} \text{ , } \varepsilon = -\frac{d\Phi_B}{dt} \\
&\text{跨科補充：} E=hf=\frac{hc}{\lambda} \text{ , } E=mc^2 \text{ , } N_t=N_0 (\frac{1}{2})^n
[\end{aligned}\]

\end{document}
